\begin{proposition}
\label{proposition-characterize-locally-finite-presentation}
Let $S$ be a scheme. Let $f : X \to Y$ be a morphism of algebraic
spaces over $S$. The following are equivalent:
\begin{enumerate}
\item The morphism $f$ is a morphism of algebraic spaces which is
locally of finite presentation, see
Morphisms of Spaces,
Definition \ref{spaces-morphisms-definition-locally-finite-presentation}.
\item The morphism $f : X \to Y$ is limit preserving as
a transformation of functors, see
Definition \ref{definition-locally-finite-presentation}.
\end{enumerate}
\end{proposition}

\begin{proof}
Assume (1). Let $T$ be a scheme and let $y \in Y(T)$. We have to show that
$T \times_Y X$ is limit preserving over $T$ in the sense of
Definition \ref{definition-locally-finite-presentation}.
Hence we are reduced to proving that if $X$ is an algebraic space which
is locally of finite presentation over $S$ as an algebraic space, then it
is limit preserving as a functor
$X : (\Sch/S)_{fppf}^{opp} \to \textit{Sets}$.
To see this choose a presentation $X = U/R$, see
Spaces, Definition \ref{spaces-definition-presentation}.
It follows from
Morphisms of Spaces,
Definition \ref{spaces-morphisms-definition-locally-finite-presentation}
that both $U$ and $R$ are schemes which are locally of finite presentation
over $S$. Hence by
Limits, Proposition
\ref{limits-proposition-characterize-locally-finite-presentation}
we have
$$
U(T) = \colim U(T_i), \quad
R(T) = \colim R(T_i)
$$
whenever $T = \lim_i T_i$ in $(\Sch/S)_{fppf}$. It follows
that the presheaf
$$
(\Sch/S)_{fppf}^{opp} \longrightarrow \textit{Sets}, \quad
W \longmapsto U(W)/R(W)
$$
is limit preserving. Hence by
Lemma \ref{lemma-sheafify-finite-presentation}
its sheafification $X = U/R$ is limit preserving too.

\medskip\noindent
Assume (2). Choose a scheme $V$ and a surjective \'etale morphism
$V \to Y$. Next, choose a scheme $U$ and a surjective \'etale morphism
$U \to V \times_Y X$. By
Lemma \ref{lemma-base-change-locally-finite-presentation}
the transformation of functors $V \times_Y X \to V$ is limit preserving. By
Morphisms of Spaces,
Lemma \ref{spaces-morphisms-lemma-etale-locally-finite-presentation}
the morphism of algebraic spaces $U \to V \times_Y X$ is locally
of finite presentation, hence limit preserving as
a transformation of functors by the first part of the proof. By
Lemma \ref{lemma-composition-locally-finite-presentation}
the composition $U \to V \times_Y X \to V$ is limit preserving
as a transformation of functors. Hence
the morphism of schemes $U \to V$ is locally of finite presentation by
Limits, Proposition
\ref{limits-proposition-characterize-locally-finite-presentation}
(modulo a set theoretic remark, see last paragraph of the proof).
This means, by definition, that (1) holds.

\medskip\noindent
Set theoretic remark. Let $U \to V$ be a morphism of
$(\Sch/S)_{fppf}$. In the statement of
Limits, Proposition
\ref{limits-proposition-characterize-locally-finite-presentation}
we characterize $U \to V$ as being locally of finite presentation
if for {\it all} directed inverse systems $(T_i, f_{ii'})$ of affine schemes
over $V$ we have $U(T) = \colim V(T_i)$, but in the current setting
we may only consider affine schemes $T_i$ over $V$ which are (isomorphic to)
an object of $(\Sch/S)_{fppf}$. So we have to make sure that there
are enough affines in $(\Sch/S)_{fppf}$ to make the proof work.
Inspecting the proof of (2) $\Rightarrow$ (1) of
Limits, Proposition
\ref{limits-proposition-characterize-locally-finite-presentation}
we see that the question reduces to the case that $U$ and $V$ are affine.
Say $U = \Spec(A)$ and $V = \Spec(B)$. By construction
of $(\Sch/S)_{fppf}$ the spectrum of any ring of cardinality
$\leq |B|$ is isomorphic to an object of $(\Sch/S)_{fppf}$.
Hence it suffices to observe that in the "only if" part of the proof of
Algebra, Lemma \ref{algebra-lemma-characterize-finite-presentation}
only $A$-algebras of cardinality $\leq |B|$ are used.
\end{proof}
